\section{Discussion}

\subsection{Interpretation of Results}

The StageBridge framework was developed to test the niche gating hypothesis: that local epithelial-stromal-immune neighborhood structure modulates cell-state transitions between LUAD initiation stages. The experimental results provided evidence relevant to both formal hypotheses articulated in the introduction.

\textbf{[INTERPRETATION: Discuss whether H1 (receiver-centered conditioning improves representation quality) was supported by the ablation results. Reference specific metrics from the ablation table comparing full model vs.\ symmetric/no-niche variants. If conditioned transitions significantly outperformed unconditioned transitions, this supports the niche gating hypothesis.]}

\textbf{[INTERPRETATION: Discuss whether H2 (dual-reference geometry improves generalization) was supported by comparing dual-reference vs.\ single-reference ablations on held-out donors. Reference kBET scores and cross-donor performance metrics. Analyze whether HLCA-only or LuCA-only ablations showed systematic weaknesses at specific disease stages.]}

The baseline comparison implemented a hierarchy of architectural inductive biases designed to isolate specific contributions. PoolingMLP served as a minimal baseline without permutation structure. DeepSets added permutation invariance through the decomposition theorem $\rho(\sum_{x \in X} \phi(x))$ but without attention mechanisms. SetTransformer introduced attention-based interactions but without spatial organization. GraphSAGE added symmetric spatial message passing. The full StageBridge model combined receiver-centered asymmetric attention with dual-reference anchoring and explicit spatial ring discretization.

\textbf{[INTERPRETATION: Discuss which architectural transitions produced the largest performance gains. Was the transition from GraphSAGE to StageBridge (adding receiver-centering and dual references) larger than the transition from SetTransformer to GraphSAGE (adding spatial structure)? What does this reveal about the relative importance of asymmetric niche encoding versus symmetric neighborhood aggregation?]}

The continuous flow matching transition model learned vector fields that transported cell state distributions between adjacent disease stages. The use of Sinkhorn-regularized optimal transport couplings to define pseudo-pairings addressed the fundamental identifiability challenge inherent to cross-sectional data, where true cell-to-cell correspondences across time were unavailable. Each transition in the LUAD progression ladder (Normal $\to$ AAH, AAH $\to$ AIS, AIS $\to$ MIA, MIA $\to$ LUAD) corresponded to biologically distinct processes: initiating mutations driving focal hyperplasia, transition to in-situ carcinoma with tissue reorganization, onset of invasion with stromal activation, and established invasion with increasing tumor heterogeneity.

\textbf{[INTERPRETATION: Discuss Sinkhorn divergence values and path straightness metrics for each transition. Were certain transitions (e.g., AIS $\to$ MIA at the invasion boundary) more difficult to model than others? Did pathway activation patterns along predicted trajectories align with known LUAD progression biology, such as EMT signatures at the invasion transition?]}


\subsection{Architectural Insights}

The hierarchical Set Transformer design with spatial ring discretization encoded a specific hypothesis about how cells integrate information from their microenvironment. By grouping neighbors into concentric shells at radii $[0, 50)$, $[50, 100)$, $[100, 150)$, and $[150, 200)$ micrometers, and applying separate ISAB modules within each ring before cross-ring attention, the architecture assumed that distance-dependent information aggregation was beneficial. This design reflected biological intuition that immediate neighbors (contact-dependent signaling) convey different information than more distant cells (paracrine signaling, tissue organization).

\textbf{[INTERPRETATION: Describe the learned distance-dependent attention patterns. Did ring 1 (0--50$\mu$m) receive consistently higher attention weights than outer rings? Were there cell-type-specific patterns where certain receivers preferentially attended to specific sender populations? Did attention patterns change systematically across disease stages, potentially reflecting progression-dependent changes in niche communication?]}

The receiver-centered formulation differed fundamentally from symmetric message-passing approaches by explicitly designating one cell as the prediction target while treating neighbors as information sources. This asymmetry reflected the biological hypothesis that a cell's transcriptional state was shaped by extrinsic microenvironmental signals in a directional manner---the niche influenced the receiver cell's state and fate, rather than bidirectional co-adaptation occurring within the timescale captured by cross-sectional snapshots. The ablation comparing receiver-centered attention against symmetric GraphSAGE directly tested whether this asymmetry was beneficial.

\textbf{[INTERPRETATION: Compare attention entropy and weight distributions between StageBridge and the symmetric GraphSAGE ablation. Did asymmetric receiver-centered attention learn qualitatively different neighborhood weighting? Was the receiver embedding more informative for downstream stage discrimination when computed via asymmetric versus symmetric aggregation?]}

The dual-reference geometry anchored query cells to complementary biological coordinate systems through HLCA (healthy lung) and LuCA (lung cancer) embeddings. The L2 normalization before concatenation prevented scale dominance, while percentile-rank calibration of k-NN distances ensured that confidence scores were comparable despite differences in atlas density. Post-hoc temperature scaling further improved confidence reliability by minimizing Expected Calibration Error (ECE), ensuring that stated confidence levels corresponded to empirical accuracy rates. This design enabled the model to simultaneously assess each cell's proximity to healthy states and similarity to known malignant phenotypes.

Beyond simple concatenation, we evaluated learned and confidence-weighted fusion strategies. Learned fusion allowed the model to discover optimal reference weighting during training, while confidence-weighted fusion dynamically adjusted weights per cell based on mapping quality. These alternatives address the question of whether fixed concatenation is optimal or whether adaptive weighting better handles transitional phenotypes poorly captured by either atlas alone.

\textbf{[INTERPRETATION: Analyze the distribution of HLCA versus LuCA confidence scores across disease stages. Normal cells should show higher HLCA confidence; LUAD cells should show higher LuCA confidence. Were intermediate stages characterized by intermediate confidence in both references, bimodal distributions, or population-specific patterns where some cells retained healthy-like profiles while others showed malignant similarity? Report the learned fusion weight at convergence---did the model favor HLCA or LuCA, and did this vary by disease stage?]}


\subsection{Comparison to Prior Work}

The StageBridge approach shared motivations with several recent methods while differing in specific architectural choices. Nicheformer \cite{Tejada-Lapuerta2025-wz} demonstrated that spatial context improved cell representations through a transformer trained on both dissociated and spatial transcriptomics, but operated on fixed neighborhood definitions without explicit receiver-centered asymmetry or dual-reference anchoring. CellFlow \cite{Klein2025-sr} applied flow matching to single-cell phenotype prediction with strong results on perturbation response, but focused on chemical and genetic perturbations rather than disease progression and did not condition flows on spatial niche context. CellOT \cite{Bunne2023-bm} used neural optimal transport for perturbation modeling with demonstrated generalization to held-out patients, but learned pointwise transport maps rather than continuous vector fields and did not incorporate spatial structure.

The COVET representation \cite{Haviv2025-kd} provided a complementary approach to encoding spatial context through gene-gene covariance structure across cells in the niche, defining OT-based distance metrics between niche representations. StageBridge differed by using explicit cell-type composition and expression statistics within spatial rings rather than covariance structure, and by integrating niche context into a continuous generative model of disease progression rather than a conditional VAE for spatial imputation.

\textbf{[INTERPRETATION: Discuss how StageBridge performance compared to expectations based on prior methods. Were there specific metrics or transitions where the architectural innovations produced larger or smaller gains than anticipated based on the literature?]}


\subsection{Limitations}

Several limitations of the present work warrant discussion.

\paragraph{Spatial mapping uncertainty.} The spatial mapping step that assigned tissue coordinates to snRNA-seq cells introduced uncertainty that propagated to downstream neighborhood construction. While multiple spatial backends were benchmarked and the best-performing method was selected, spatial assignment accuracy varied across cell types and tissue regions. Cell types with distinctive marker profiles that were well-represented in the spatial reference showed higher mapping confidence than rare or transcriptionally ambiguous populations. Regions of high cellular density or complex tissue architecture posed particular challenges for accurate localization.

\textbf{[LIMITATION: Report spatial mapping accuracy stratified by cell type and disease stage. Which populations showed lowest confidence? How did neighborhood construction errors affect learned niche representations for those populations?]}

\paragraph{Cross-sectional identifiability.} The cross-sectional nature of the data precluded direct validation of learned transition dynamics against true longitudinal trajectories. While flow matching produced vector fields that transported source distributions to target distributions with quantifiably low Sinkhorn divergence, individual cell trajectories represented statistical inference rather than observed biological paths. The optimal transport coupling provided an alignment prior that assumed cells transitioned to similar transcriptional states, but this reflected statistical parsimony rather than verified biological ancestry. Cells that underwent dramatic transcriptional reprogramming during transitions might be systematically mispaired.

\paragraph{Trajectory validation limitations.} Without lineage tracing or longitudinal sampling of the same tissue, the biological plausibility of learned trajectories could only be assessed indirectly through pathway activity analysis and marker gene dynamics. Trajectories that appeared biologically sensible based on known progression signatures might nonetheless reflect shortcuts through transcriptional space that real cells never traverse.

\paragraph{Reference atlas coverage.} The dual-reference geometry was limited to two atlases (HLCA and LuCA), which may not have captured all relevant biological variation in the LUAD progression cohort. Cells with phenotypes poorly represented in either reference---such as novel transitional states, rare stromal subtypes, or patient-specific populations---received low confidence scores in both coordinate systems, potentially degrading their representation quality and downstream utility.

\textbf{[LIMITATION: Report the fraction of cells with confidence below threshold in both references. Were there specific cell types, disease stages, or patients where dual-reference anchoring provided limited benefit? Did low-confidence cells cluster together in embedding space, suggesting systematic gaps in reference coverage?]}

\paragraph{Geometric assumptions.} The ring-based spatial discretization imposed a particular geometric assumption about niche organization, treating the microenvironment as concentrically organized around each receiver cell. While this parameterization was biologically motivated by the radial decay of paracrine signaling and contact-dependent communication, tissue architecture in solid tumors is often anisotropic. Stromal barriers, vascular networks, and lobular organization create directional biases that concentric rings cannot capture. Alternative parameterizations---sector-based partitioning, adaptive distance thresholds, or learned spatial attention without explicit discretization---might better model complex tissue geometry.

\paragraph{Single timepoint per stage.} Each disease stage in the LUAD cohort represented a single cross-sectional snapshot, conflating biological variation within the stage with technical and sampling variation. The model could not distinguish whether cells at the boundary between stages represented true transitional states versus measurement noise. Future cohorts with multiple independent samples per stage would enable more robust characterization of within-stage heterogeneity.

\paragraph{Computational constraints.} The scale of hyperparameter exploration was limited by computational requirements. Self-supervised pretraining and CFM-OT training each required substantial GPU hours, constraining the breadth of architectural variants evaluated. The reported configuration represented a reasonable operating point but not necessarily the global optimum.

\paragraph{Addressed limitations.} Several potential limitations identified during development were addressed through implementation improvements. First, the fusion strategy between HLCA and LuCA embeddings was extended beyond simple concatenation to include learned weighted fusion (where the optimal weight is discovered during training) and confidence-weighted fusion (where mapping quality dynamically adjusts per-cell weighting). Ablation studies quantify the impact of these alternatives. Second, confidence calibration via temperature scaling was implemented to improve the reliability of mapping confidence scores, reducing Expected Calibration Error and ensuring that confidence values better reflect true accuracy. Third, SSL loss weights were made configurable with systematic ablations comparing the default 70\% masked reconstruction emphasis against equal weighting and masked-only variants. These improvements transform what would otherwise be architectural assumptions into empirically validated design choices.


\subsection{Future Directions}

The StageBridge framework established foundations for several extensions that address current limitations and expand the scope of biological questions accessible to transition modeling.

\paragraph{Non-Euclidean latent geometry.} The current dual-reference geometry employed Euclidean embeddings with simple concatenation. However, cell type hierarchies exhibit natural tree structure that Euclidean space represents inefficiently, while disease progression involves branching trajectories where some paths lead to distinct outcomes. Future work will explore hyperbolic embeddings for hierarchical cell type relationships, where distances grow exponentially with tree depth, and spherical geometry for periodic processes such as cell cycle. Product manifolds combining Euclidean, hyperbolic, and spherical components could capture complementary biological structure within a unified framework \cite{Palma2025-uk}.

\paragraph{Neural stochastic differential equations.} While flow matching provided efficient training of continuous dynamics, the learned vector field was deterministic given state and time. Biological state transitions involve intrinsic stochasticity from transcriptional bursting, epigenetic fluctuations, and microenvironmental variability. Neural SDE backends would model state-dependent diffusion, enabling uncertainty quantification that reflects biological noise rather than only epistemic uncertainty from limited training data. Score matching objectives could train such models while maintaining computational tractability.

\paragraph{Phase portrait and attractor analysis.} The learned vector field implicitly defined a dynamical system on the cell state manifold. Future work will extract explicit phase portraits from this system, identifying fixed points (stable cell states), basins of attraction (cell fate commitments), and separatrices (decision boundaries between fates). For LUAD progression, this analysis could reveal whether the AIS-to-MIA invasion transition represents crossing a separatrix between non-invasive and invasive attractors, and whether niche composition modulates the position of that boundary.

\paragraph{Cohort transport layer.} The current model operated at the cell and neighborhood level. A natural extension would embed entire patient-stage samples as point clouds in a transport-aware space, where distances between samples reflected optimal transport cost between their cell populations. This cohort transport layer would enable patient-level trajectory analysis, identification of outlier progression patterns, and stratification of patients by their position in a learned progression geometry.

\paragraph{Destination-conditioned transitions.} The LUAD progression ladder represents within-organ evolution, but metastasis involves cross-organ transitions where destination tissue shapes arriving cell phenotypes. Future versions will support destination conditioning, where the transition model receives not only source niche context but also target tissue identity. This extension would enable questions about metastatic tropism---why certain primary tumors preferentially seed specific distant sites---and plasticity---how cells adapt their transcriptional programs to foreign microenvironments. Lung to brain metastasis represents a particularly tractable initial target given the distinct microenvironmental differences.

\paragraph{Higher-order interactions via hypergraphs.} Pairwise cell-cell interactions captured by the current graph structure may miss biologically important higher-order effects. Immune synapses involve coordinated three-way interactions between antigen-presenting cells, T cells, and target cells. Clone-niche-stage relationships involve higher-order dependencies where evolutionary state modulates how cells respond to microenvironmental signals. Hypergraph extensions would represent such interactions as hyperedges connecting multiple cells, enabling the model to learn non-additive effects that pairwise graphs cannot capture.

\paragraph{Birth-death dynamics.} The current model learned state transitions but not population dynamics. Extending to neighborhood-aware birth-death processes would enable questions about how niche composition affects not only cell state but also proliferation and survival. This connects to foundational work on tissue dynamics and would require integrating optimal transport for state transitions with branching process models for population changes.

\paragraph{Causal intervention modeling.} The learned representations and transition dynamics could support counterfactual prediction---simulating how cell states would evolve under hypothetical interventions. Given a cell's niche-conditioned embedding, the framework could potentially predict responses to therapeutic perturbations by modeling drug effects as shifts in the embedding space or modifications to the transition vector field. This connects StageBridge to the broader virtual cell paradigm where computational models support experimental design and hypothesis generation.


\subsection{Broader Impact}

The StageBridge framework addressed a fundamental challenge in computational biology: learning predictive models of disease progression from cross-sectional measurements when longitudinal tracking is infeasible. Success in this domain could accelerate understanding of cancer evolution, identify microenvironmental targets for therapeutic intervention, and inform early detection strategies by characterizing the molecular features of pre-invasive lesions.

The receiver-centered formulation emphasized that cellular phenotypes emerge from interaction between intrinsic transcriptional programs and extrinsic microenvironmental signals. This perspective aligned with growing recognition in cancer biology that the tumor microenvironment plays central roles in progression, metastasis, and treatment response. Computational methods that explicitly model niche context complement experimental approaches that manipulate microenvironmental composition, potentially identifying which stromal or immune populations represent therapeutic targets at specific disease stages.

The dual-reference geometry demonstrated a general strategy for anchoring query data to multiple biological coordinate systems. As single-cell atlases proliferate across tissues, diseases, and species, reference-based methods that integrate multiple anchors will become increasingly valuable. The confidence calibration approach ensured that distances in different reference spaces remained comparable, addressing a practical challenge that recurs whenever embeddings from different sources must be combined.

The continuous flow matching transition model provided a generative framework for simulating disease progression trajectories. While the present work focused on representation learning and trajectory inference, the learned vector fields could support counterfactual prediction in future extensions---generating hypothetical cell states under alternative progression scenarios or therapeutic interventions. Such capabilities connect to broader efforts in causal inference for single-cell biology and the emerging virtual cell paradigm that envisions computational models capable of predicting cellular responses to arbitrary perturbations.

\textbf{[IMPACT: Discuss any specific biological insights that emerged from attention analysis or trajectory modeling. Were there unexpected findings about niche composition changes, stage-specific cell-cell interactions, or gene expression dynamics that could inform follow-up experimental studies or clinical biomarker development?]}

